\documentclass[11pt]{article}
    \usepackage[utf8]{inputenc}
    \usepackage[T1]{fontenc}
    \usepackage[margin=1in]{geometry}
    \usepackage{hyperref}
    \usepackage{enumitem}
    \title{Indi & Geometry | Part 1 (Day 1)}
    \author{Piter Garcia}
    \date{March 20, 2026}
    \begin{document}
    \maketitle
    \section*{Quick Scan}
    \begin{itemize}[leftmargin=*]
    \item Grade: Grade 1
    \item Workbook sessions: Session 19
    \item Class blocks: Day 3 12:25-1:15 Gargana, Day 3 2:15-3:05 Polfleit, Day 5 12:25-1:15 Montgomery
    \item Reference file: indi-and-geometry-part-1-day-1.bib
    \end{itemize}
    \section*{School Planning Goals and Inclusion}
    \begin{itemize}[leftmargin=*]
    \item What is expected: Indi Cars | Indi Kits | Path | Paper & Pencil | Shape Challenge Cards
    \item What we achieved: No direct transcript match is attached yet. Treat this lesson as workbook-backed and enrich it when a matching classroom transcript is available.
    \item What needs improvement: stronger proof, cleaner assessment notes, and tighter lesson artifacts.
    \item How to improve: add transcript proof, one saved artifact, and clearer success criteria.
    \item Inclusion focus: use visual modeling, chunked directions, predictable routines, and flexible participation roles.
    \end{itemize}
    \section*{Official References}
    \begin{itemize}[leftmargin=*]
    \item \url{https://csteachers.org/k12standards/interactive/}
\item \url{https://dmmedia.sphero.com/email-marketing/Sphero-Edu/SpheroEdu-k12-teacher-resource-guide-v1_updated050818.pdf}

    \end{itemize}
    \end{document}
